\begin{abstract}
Receding-horizon control with a learned dynamics model is sensitive to both
forecast-origin state error and accumulated transition error. This paper
presents a Complementary-Filtered Residual World Model (CFR-WM) that applies a
mask-gated, gain-0.5 context update before free rollout with nominal dynamics
and a bounded neural residual. Across five nonlinear controlled systems, three
test regimes, and eight paired seeds, CFR-WM reduces 25-step normalized root
mean squared error by \PrimaryImprovement\% relative to a physics-guided neural
baseline (95\% interval [\PrimaryRelativeCILow\%,
\PrimaryRelativeCIHigh\%], $p=\PrimaryPValue$ against a 5\% margin). The
corrupt-regime improvement is \CorruptImprovement\%, whereas clean-regime
noninferiority is not established. In \ClosedLoopEpisodes{} paired simulated
receding-horizon control episodes, CFR-WM reduces realized normalized state
loss by \ClosedLoopImprovement\% (95\% interval [\ClosedLoopCILow\%,
\ClosedLoopCIHigh\%], $p=\ClosedLoopPValue$) relative to the same baseline.
However, filtering followed by nominal dynamics slightly outperforms the
complete model in both open- and closed-loop evaluations. On held-out measured
NanoDrone flights, CFR-WM lowers 50-step position error relative to the matched
unfiltered residual model in all four context conditions, but increases
full-state error in all eight condition--horizon comparisons; its predictive
intervals also under-cover. The results identify forecast-origin filtering as
the supported source of robustness under corrupted observations and delimit
the evidence for the learned residual and uncertainty output.
\end{abstract}
